\documentclass{article}

\usepackage{corl_2026}
\usepackage{amsmath}
\usepackage{amssymb}
\usepackage{bm}
\usepackage{booktabs}
\usepackage{array}
\usepackage{graphicx}
\usepackage{textcomp}
\usepackage[UTF8]{ctex}

\title{面向学习型固定翼机动飞行的滚动时域目标流方法}

\author{匿名作者}

\newcommand{\rhtso}{RH-TSO}
\newcommand{\vt}{v_t}

\begin{document}
\maketitle

%===============================================================================

\begin{abstract}
固定翼机动飞行不能简单化为航点跟踪。固定翼飞机的姿态、速度、升力方向、能量状态和作动器带宽强耦合，因此飞机即使在位置上接近参考轨迹，也可能以错误的姿态或能量状态完成机动。
本文研究一种组合式方法：首先训练一个冻结的、无 PID 的四元数条件强化学习飞行技能，使其根据目标航向、俯仰、滚转和空速，直接输出油门、升降舵、副翼、方向舵和减速板指令。
长时域机动被表示为该技能可执行的目标流；目标流参数由滚动时域目标流优化（Receding-Horizon Target-Stream Optimization, \rhtso）在线选择。\rhtso 使用冻结策略的短时域闭环 rollout，对候选目标流进行几何、能量、安全和平滑性代价评估。
在四类代表性机动上，最优目标流参数具有明显任务依赖性：水平曲线和温和三维曲线偏好更短的前瞻距离和更低的目标速度，而 $90^\circ$ 垂直拉起偏好更长的前瞻距离和更高的目标速度。相对于默认目标流，S 形曲线、8 字形、螺旋线和垂直拉起的 CTE-P90 分别改善 39\%、2\%、24\% 和 47\%。
能量感知微调将技能扩展到垂直弧线，但几何感知评估给出了清晰边界：$60^\circ$--$150^\circ$ 垂直弧线为 B 级，$180^\circ$ 半筋斗仍为失败案例。
\end{abstract}

\keywords{固定翼空中机器人，强化学习，四元数控制，目标流，滚动时域优化}

%===============================================================================

\section{引言}

固定翼飞机不能像多旋翼那样在航点处悬停、原地转向再继续飞行。它能否转弯、爬升或倒飞，取决于姿态如何改变升力方向、俯仰如何在速度和高度之间交换能量，以及气动舵面还剩多少控制裕度。因此，单纯跟踪几何曲线并不能完整描述固定翼机动。一条轨迹可以具有很小的横偏误差（cross-track error, CTE），但同时伴随侧滑、能量损失、环面偏转或接近失速。

经典固定翼飞行栈通常采用嵌套控制结构：规划器给出高层参考，PID 或自动驾驶仪内环跟踪姿态和速度，安全逻辑负责限制飞行包线。这类系统在设计工况内非常有效，但对于学习型机动组合存在两个限制。第一，学习模块被隔离在作动器级动力学之外，而真正决定激进指令是否可执行的是作动器和气动响应。第二，欧拉角目标在接近垂直和倒飞姿态时表达不自然，而这些正是环类机动最容易暴露问题的区域。

本文训练一个可复用的飞行技能，使其从状态和目标姿态--速度指令直接映射到作动器指令。策略观测四元数姿态误差、机体系目标方向、速度误差、飞机状态和上一时刻动作；输出油门、升降舵、副翼、方向舵和减速板。该技能训练一次后被冻结。之后，规划问题不再是“应该施加什么作动器序列”，而是“应该向该技能输入什么可执行目标流”。

目标流是本文的核心接口。参考机动被转换为局部航向、俯仰、滚转和目标空速命令，冻结策略据此执行。固定转换并不够。最新目标流扫描显示，S 形曲线、8 字形和螺旋线偏好较短前瞻距离和较低目标速度，而垂直拉起偏好相反。最优方向随机动类别改变，因此不存在一个适用于所有机动的固定目标流参数。

我们提出滚动时域目标流优化（\rhtso），在短时域内优化可执行目标流参数。每个重规划步，候选目标流都通过同一个冻结策略和动力学进行闭环 rollout，再根据几何、能量、安全和平滑性代价选择最优方案。当前实现采用确定性并行打靶格点搜索，因为实验中的目标流空间维数较低；当目标流参数化变得更高维或连续时，同一形式也可以自然扩展到 CEM/MPPI 风格的采样优化。

本文贡献如下：
\begin{enumerate}
  \item 一个无 PID 的四元数条件固定翼飞行技能，直接将目标航向、俯仰、滚转和空速映射为作动器级控制。
  \item 一个可执行目标流抽象，将长时域参考机动转换为冻结技能可跟踪的局部命令，避免静态航点或静态姿态目标的局限。
  \item \rhtso：一种在目标流参数空间中进行滚动时域优化的方法，本文以可执行技能空间中的确定性并行打靶格点搜索实现。
  \item 一个能量感知垂直扩展，通过低速、能量保持、安全、垂直进展和原始任务回放项改善大姿态垂直弧线能力，但不声称解决 $180^\circ$ 半筋斗。
  \item 一个几何感知评估协议，说明为什么 CTE 单独不足以评价环类固定翼机动，并给出当前能力边界。
\end{enumerate}

所有实验均在含六自由度 F-16 气动模型的高保真固定翼仿真中完成。该飞机模型是实验载体，不是本文的主要贡献。

%===============================================================================

\section{相关工作}

\paragraph{敏捷飞行中的学习与控制接口。}
多项 CoRL 工作表明，将学习模块放在可检查的中间接口上，往往比替代整套自治系统更稳健。Deep Drone Racing 将机载图像映射为航点和期望速度，再交给规划器和控制器执行~\citep{kaufmann2018deep}。基于学习的航点导航也将学习感知与模型规划、反馈控制相结合~\citep{bansal2019combining}。这些工作启发了本文的接口设计，但固定翼机动所需的可执行对象不同：它不仅需要局部位置目标，还需要完整姿态和速度命令。

\paragraph{可行性约束下的轨迹跟踪。}
DATT 将四旋翼轨迹跟踪表述为受作动器限制、扰动和不可行参考影响的问题~\citep{huang2023datt}。用于四旋翼的学习式时间分配方法也强调，生成的轨迹必须通过可飞性检验~\citep{ryou2022real}。本文面对类似的可行性问题，但约束由冻结的固定翼技能闭环体现。\rhtso 不优化作动器命令，也不优化多项式轨迹；它优化的是学习技能能够执行的目标流。

\paragraph{直接学习控制。}
端到端四旋翼集群控制表明，策略可以从局部观测直接输出推力命令~\citep{batra2021decentralized}。当传统控制器设计困难时，学习动力学和控制也被用于软体多旋翼等系统~\citep{deng2020soft}。本文的技能同样是作动器级直接控制，但仍以显式目标流为条件。这一设计保留了接口的可检查性，同时避免手工调节内环。

\paragraph{仿真基础设施与评估。}
Flightmare 展示了速度、保真度和学习吞吐之间可调仿真基础设施的重要性~\citep{song2020flightmare}。本文同样依赖高吞吐仿真，但研究对象不是仿真器本身。我们的评估问题更具体：固定翼机动不能只看位置误差，还需要联合姿态、切线、翼面、能量和安全指标。

%===============================================================================

\section{方法}

\begin{figure}[t]
  \centering
  \fbox{\begin{minipage}{0.94\linewidth}
  \small
  \textbf{系统概览。}
  参考机动 $\rightarrow$ 目标流生成器 / \rhtso\ $\rightarrow$
  局部 $(\psi,\theta,\phi,\vt^{\star})$ 目标流 $\rightarrow$
  冻结四元数条件飞行技能 $\pi_\theta$ $\rightarrow$
  作动器指令 $(T,\delta_e,\delta_a,\delta_r,\delta_b)$ $\rightarrow$
  固定翼动力学 $\rightarrow$ 几何感知评估。
  \end{minipage}}
  \caption{主要接口。\rhtso 不在作动器空间中搜索，而是在目标流参数空间中搜索，并通过冻结技能的闭环 rollout 评估每个候选方案。}
  \label{fig:system-cn}
\end{figure}

\subsection{四元数条件作动器级直接飞行技能}

底层技能是一个通过 PPO 训练的循环 actor-critic 策略，训练后在组合实验中保持冻结。该策略不使用 PID：它直接输出仿真器的五个作动器通道。表~\ref{tab:policy-interface-cn} 总结了接口。

\begin{table}[t]
  \caption{飞行技能接口。目标是局部姿态和空速，而不是航点。}
  \label{tab:policy-interface-cn}
  \centering
  \begin{tabular}{p{0.24\linewidth}p{0.64\linewidth}}
    \toprule
    \textbf{组成} & \textbf{变量} \\
    \midrule
    目标命令 & 目标航向 $\psi^\star$、俯仰 $\theta^\star$、滚转 $\phi^\star$、目标空速 $\vt^\star$ \\
    姿态观测 & 最短旋转四元数误差 $q_e$ 的向量部 \\
    方向与速度 & 机体系目标方向、归一化速度误差、归一化速度、归一化高度 \\
    飞机状态 & 角速率、攻角/侧滑特征、上一时刻动作 \\
    动作输出 & 油门、升降舵、副翼、方向舵、减速板 \\
    \bottomrule
  \end{tabular}
\end{table}

设 $q$ 为当前飞机姿态，$q^\star$ 为由目标航向、俯仰、滚转构造的目标姿态。策略观测最短旋转误差
\begin{equation}
  q_e = q^\star \otimes q^{-1}, \qquad q_e \leftarrow \operatorname{sign}(q_{e,w})q_e ,
  \label{eq:quat-error-cn}
\end{equation}
并使用 $q_e$ 的向量部作为有界三维姿态误差信号。对应的测地距离为
\begin{equation}
  d_{\mathrm{SO(3)}}(q,q^\star) = 2\arccos\left(|\langle q,q^\star\rangle|\right).
\end{equation}
我们采用四元数目标，是为了避免垂直和倒飞姿态附近的表示不连续。除非补充并核验欧拉基线消融表，本文不作“四元数显著优于欧拉角”的定量结论。

基准奖励结合姿态跟踪、速度跟踪、存活和高度/安全项。关键设计不是 PPO 本身，而是可复用的技能接口：策略在没有 PID 或自动驾驶仪内环的情况下，学习对动态姿态--速度目标进行作动器级跟踪。

\subsection{可执行目标流组合}

冻结技能接收局部目标流
\begin{equation}
  \tau_k = (\psi_k^\star,\theta_k^\star,\phi_k^\star,\vt_k^\star).
\end{equation}
给定参考机动 $\mathcal{M}$ 和目标流参数 $\lambda$，目标流生成器 $G$ 产生
\begin{equation}
  \tau_{t:t+H} = G(\mathcal{M}, x_t, \lambda).
\end{equation}
实验中，$\lambda$ 包括前瞻距离、目标空速，以及用于垂直塑形的俯仰偏移。这些参数高于作动器层级，但低于符号化机动标签；它们决定参考几何的哪一部分以怎样的局部命令形式交给飞行技能执行。

静态航点不适合作为这一问题的接口，因为单个点无法指定升力方向、姿态准备或能量管理。静态姿态--速度目标也不够，因为它无法传达期望姿态沿机动如何变化。目标流提供了参考几何与作动器级学习技能之间的动态局部契约。

\subsection{滚动时域目标流优化}

\begin{figure}[t]
  \centering
  \fbox{\begin{minipage}{0.94\linewidth}
  \small
  \textbf{目标流优化。}
  在状态 $x_t$，枚举或采样候选目标流参数 $\lambda_i$。
  对每个候选：生成 $\tau_{t:t+H}^{(i)}$，闭环 rollout
  $u_k=\pi_\theta(x_k,\tau_k^{(i)})$、$x_{k+1}=f(x_k,u_k)$，计算 $J^{(i)}$。
  选择最低代价目标流，执行第一段，然后重规划。
  \end{minipage}}
  \caption{\rhtso 通过冻结策略的闭环执行评估候选目标流，而不是只按开环几何打分。}
  \label{fig:rhtso-cn}
\end{figure}

\rhtso 求解短时域目标流参数优化问题：
\begin{align}
  \lambda_t^\star
  &= \arg\min_{\lambda \in \Lambda} J(x_{t:t+H}^{\lambda},u_{t:t+H-1}^{\lambda},\mathcal{M}), \label{eq:rhtso-cn}\\
  \tau_k^\lambda &= G(\mathcal{M},x_k^\lambda,\lambda), \\
  u_k^\lambda &= \pi_\theta(x_k^\lambda,\tau_k^\lambda), \\
  x_{k+1}^\lambda &= f(x_k^\lambda,u_k^\lambda),
  \qquad k=t,\ldots,t+H-1 .
\end{align}
优化变量是目标流参数，而不是作动器命令。因此搜索维度较低，每个候选方案也具有物理解释。

当前实现中的代价包含四项：
\begin{equation}
  J = w_gJ_{\mathrm{geo}} + w_eJ_{\mathrm{energy}} + w_sJ_{\mathrm{safety}} + w_uJ_{\mathrm{smooth}}.
  \label{eq:cost-cn}
\end{equation}
$J_{\mathrm{geo}}$ 衡量到参考轨迹的距离以及与参考切线、机动平面的对齐；$J_{\mathrm{energy}}$ 惩罚相对机动需求的比能损失；$J_{\mathrm{safety}}$ 惩罚低速、过大攻角或侧滑、超 G 和高度违规；$J_{\mathrm{smooth}}$ 惩罚作动器变化。

我们采用确定性并行打靶格点搜索，是因为当前实验中的目标流空间维数较低，该方法简单且在给定格点上最优。当目标流空间变得更高维或连续时，同一形式可以自然扩展到 CEM/MPPI 风格的采样优化。

\subsection{能量感知垂直扩展}

基准技能在水平和中等姿态目标上较可靠，但垂直弧线会暴露能量损失和环面漂移问题。我们从基准四元数技能出发进行微调，在加入垂直弧线目标的同时回放原始任务以保持水平飞行能力。奖励中加入低速惩罚、能量保持项、高度和垂直进展塑形、攻角/侧滑/G 安全惩罚以及动作平滑项。环类阶段还包含速度切线、机头切线、机头速度和翼面对齐等几何项。

这一扩展应被理解为能力边界的推进，而不是完整特技飞行的解决方案。它改善了 $60^\circ$--$150^\circ$ 垂直弧线，但 $180^\circ$ 半筋斗仍不在当前可靠包线内。

\subsection{几何感知机动指标}

CTE 是必要指标，但不充分。对于环类机动，我们还评估：
\begin{align}
  e_v &= \arccos(\hat v \cdot \hat t_{\mathrm{ref}}), &
  e_n &= \arccos(\hat n_{\mathrm{body}} \cdot \hat t_{\mathrm{ref}}), \\
  e_{nv} &= \arccos(\hat n_{\mathrm{body}} \cdot \hat v), &
  e_w &= \arccos(|\hat w_{\mathrm{body}} \cdot \hat n_{\mathrm{plane}}|), \\
  e_q &= d_{\mathrm{SO(3)}}(q,q_{\mathrm{ref}}).
\end{align}
其中 $e_v$ 为速度切线误差，$e_n$ 为机头切线误差，$e_{nv}$ 为机头速度误差，$e_w$ 为翼面误差，$e_q$ 为四元数测地误差。这些指标与空速、攻角/侧滑、过载、高度和能量包线一起报告。

%===============================================================================

\section{实验设置}

仿真器使用六自由度固定翼动力学模型、气动查表、作动器限制和 50 Hz 物理积分。RL 策略以较低控制频率调用，输出归一化作动器指令，再解码到仿真通道。

目标流实验使用冻结的能量感知技能。评估任务包括 S 形曲线、8 字形、螺旋线/温和三维曲线和 $90^\circ$ 垂直拉起。默认目标流为 $(L,\vt^\star)=(1000~\mathrm{m},250~\mathrm{m/s})$。报告的格点结果比较该默认设置与同一闭环评估下的任务自适应选择。

环质量评估覆盖 $60^\circ$ 到 $180^\circ$ 的垂直弧线。当前稿件将详细数值表列为待核验项；本文只保留已诊断的定性边界：$60^\circ$--$150^\circ$ 弧线为 B 级，$180^\circ$ 半筋斗为 Fail。

%===============================================================================

\section{实验结果}

\subsection{底层技能跟踪}

作动器级直接飞行技能能够在没有 PID 内环的情况下跟踪动态姿态--速度目标，因此可作为目标流的执行器。最终稿应加入一个紧凑表格，报告保留目标序列上的姿态测地误差、速度误差、存活率和安全违规率。由于稿件目录中没有相应评估日志，这里不复述未核验数值。

\subsection{欧拉目标与四元数目标}

采用四元数条件的原因首先是表示层面的：式~\eqref{eq:quat-error-cn} 中的误差在垂直姿态附近和倒飞阶段仍然定义良好。当前稿件只使用这一动机。只有在欧拉基线表补充并核验后，才应加入四元数相对欧拉的定量优势结论。

\subsection{目标流消融}

静态空间目标和静态姿态--速度目标都不是长时域固定翼机动的可执行描述。目标流提供时变局部命令，同时传达期望姿态以及该姿态沿参考几何如何变化。该消融需要在补齐代表性轨迹图后进行可视化说明。

\subsection{滚动时域目标流选择}

\rhtso 当前最有力的证据是目标流参数选择的任务依赖性。表~\ref{tab:rhtso-results-cn} 报告最新的二维目标流切片扫描。水平和温和三维机动偏好更短前瞻和更低目标速度；垂直拉起偏好更长前瞻和更高目标速度。因此，目标流参数不能对所有机动固定不变。

\begin{table}[t]
  \caption{最新闭环扫描中的目标流选择。最佳值与默认目标流 $(1000~\mathrm{m},250~\mathrm{m/s})$ 比较。CTE-P90 单位为米。}
  \label{tab:rhtso-results-cn}
  \centering
  \begin{tabular}{lcccc}
    \toprule
    \textbf{机动} & \textbf{最佳 $(L,\vt^\star)$} & \textbf{最佳 CTE-P90} & \textbf{默认 CTE-P90} & \textbf{改善} \\
    \midrule
    S 形曲线 & $(600,220)$ & 918 & 1503 & 39\% \\
    8 字形 & $(600,220)$ & 1017 & 1039 & 2\% \\
    螺旋线 / 温和 3D & $(600,220)$ & 524 & 691 & 24\% \\
    $90^\circ$ 垂直拉起 & $(1500,280)$ & 51 & 96 & 47\% \\
    \bottomrule
  \end{tabular}
\end{table}

这一结果并不是说 $(600,220)$ 普遍最好。它适合水平和温和三维曲线，但不适合垂直拉起。反过来，垂直机动的选择 $(1500,280)$ 对曲线任务也不是好选择。这正是需要在线优化可执行目标流参数的实验证据。

\subsection{能量感知垂直扩展}

能量感知微调通过显式塑形低速行为、能量保持、垂直进展、安全裕度和原始任务回放，将基准技能推向更大姿态的垂直弧线。当前证据应表述为能力边界，而不是已解决环飞。该技能在 $60^\circ$--$150^\circ$ 垂直弧线下达到 B 级几何质量，而 $180^\circ$ 半筋斗仍为 Fail。

\subsection{几何感知评估与能力边界}

\begin{figure}[t]
  \centering
  \fbox{\begin{minipage}{0.94\linewidth}
  \small
  \textbf{代表性机动图占位。}
  后续替换为 S 形曲线、8 字形、螺旋线和垂直拉起的并排图：
  参考轨迹、实际轨迹、选择的目标流参数以及 CTE 时间曲线。
  \end{minipage}}
  \caption{投稿前需要补齐的代表性机动图。}
  \label{fig:maneuvers-cn}
\end{figure}

\begin{figure}[t]
  \centering
  \fbox{\begin{minipage}{0.94\linewidth}
  \small
  \textbf{垂直弧线边界图占位。}
  后续替换为 $60^\circ$--$180^\circ$ 弧线的分相位曲线，对比 CTE、速度切线误差、机头切线误差、机头速度误差、翼面误差和四元数测地误差。
  \end{minipage}}
  \caption{CTE 单独评估会掩盖错误环面几何；核验后的图应显示 $60^\circ$--$150^\circ$ 的 B 级边界和 $180^\circ$ 失败。}
  \label{fig:vertical-frontier-cn}
\end{figure}

表~\ref{tab:loop-frontier-cn} 记录当前能力边界，同时不编造缺失指标。详细数值需由 loop-quality 日志补齐。

\begin{table}[t]
  \caption{环质量边界。指标数值保留为待核验占位；评级来自当前严格诊断。}
  \label{tab:loop-frontier-cn}
  \centering
  \begin{tabular}{lccccc}
    \toprule
    \textbf{弧线角度} & \textbf{CTE} & \textbf{速度切线} & \textbf{机头切线} & \textbf{翼面} & \textbf{评级} \\
    \midrule
    $60^\circ$ & 待核验 & 待核验 & 待核验 & 待核验 & B \\
    $90^\circ$ & 待核验 & 待核验 & 待核验 & 待核验 & B \\
    $120^\circ$ & 待核验 & 待核验 & 待核验 & 待核验 & B \\
    $150^\circ$ & 待核验 & 待核验 & 待核验 & 待核验 & B \\
    $180^\circ$ & 待核验 & 待核验 & 待核验 & 待核验 & Fail \\
    \bottomrule
  \end{tabular}
\end{table}

这里的评估结论是定性的，但很关键：CTE 单独可能错误排序环类机动。靠近参考轨迹并不代表翼面、机头方向、速度切线或能量状态正确。几何感知指标揭示了当前单一冻结技能系统的真实边界。

%===============================================================================

\section{讨论与局限}

\paragraph{为什么优化目标流。}
目标流空间维数低、可解释且可执行。维数低，是因为少数参数即可改变一段时域内的局部命令；可解释，是因为所选前瞻距离和速度反映了技能如何为机动做准备；可执行，是因为每个候选都通过冻结策略和动力学闭环评估后才被选择。

\paragraph{\rhtso 不声称什么。}
确定性格点搜索只在给定候选集合上具有格点最优性。它不是全局最优，不是作动器空间 MPC，也不是当前实验中格点方法优于 CEM 或 MPPI 的证据。

\paragraph{仿真范围。}
所有结果均来自仿真。真实部署还需要处理状态估计噪声、风、作动器动态、机体模型差异以及额外安全监控，这些不在本文范围内。

\paragraph{能力边界。}
当前系统不是完整特技飞行方案。水平和温和三维机动在测试目标流下较可靠；$60^\circ$--$150^\circ$ 垂直弧线在环质量评估下为 B 级；$180^\circ$ 半筋斗是失败边界。一个初步的相位门控 residual specialist 可能有助于倒飞过渡区，但本文主评估聚焦于单一冻结技能的可复用目标流组合，因此不将该 specialist 纳入主结果。

%===============================================================================

\section{结论}

本文提出了一个面向固定翼机动飞行的组合式框架，其核心是可复用的无 PID 四元数条件飞行技能和可执行目标流。\rhtso 通过冻结技能的短时域闭环 rollout 选择目标流参数，揭示了固定目标流难以跨机动类别工作的原因：水平曲线任务偏好短前瞻和较低速度，而垂直拉起偏好长前瞻和较高速度。同一评估也明确了当前能力边界。能量感知微调改善了垂直弧线，但几何感知指标表明，$180^\circ$ 半筋斗仍为 Fail，而不是已解决问题。下一步应补齐待核验的环质量表和代表性轨迹图。

%===============================================================================

\bibliography{example}

\end{document}
